\title{Byte Latent Transformer: Patches Scale Better Than Tokens}

\begin{abstract}
We present the Byte Latent Transformer ({\textsc BLT}), a novel byte-level large language model architecture that, for the first time, achieves performance on par with tokenization-based LLMs at scale, while offering substantial gains in inference speed and robustness. 
{\textsc BLT} transforms bytes into variable-length patches, which function as the main computational units within the model. These patches are defined according to the entropy of upcoming bytes, allowing the system to allocate greater computational resources and representational power in regions of high data complexity. We provide the first scaling analysis of byte-level models under controlled computational budgets, scaling up to 8 billion parameters and training on 4 trillion bytes. The findings highlight that scaling models directly on byte inputs, without reliance on fixed vocabularies, is entirely practical. The approach enhances both training and inference throughput by adaptively using longer patches in less complex contexts, alongside observed qualitative benefits in reasoning and generalization to less common cases. Taken together, for equal inference budgets, {\textsc BLT} demonstrates a markedly improved scaling curve relative to tokenized architectures by enabling simultaneous increases in both patch length and model size.
\end{abstract}

\section{Introduction}
\label{section:intro}

We present the Byte Latent Transformer~(\textbf{BLT}), an architecture that operates directly on raw byte sequences without relying on any form of tokenization. BLT, for the first time, achieves parity in performance with leading models that utilize token-based pre-processing at scale, while also delivering notable gains in both robustness and computational efficiency (\S\ref{sec:robustness}).
Current large language models (llms) are primarily end-to-end systems, yet they universally depend on tokenization—a pre-defined and heuristic approach that segments bytes into a predetermined vocabulary of tokens. This token-based step introduces compression biases based on token definition, which results in limitations such as sensitivity to particular domains or modalities~\citep{dagangetting}, decreased robustness to noisy inputs~(\S\ref{sec:robustness}), insufficient understanding of orthographic variation~\citep{edman2024cute}, and imbalances in multilingual processing capabilities~\citep{liang2023xlm,petrov2024language,limisiewicz2024myte}.

Historically, tokenization has been a crucial step because training llms on raw bytes is infeasible at large scales, given the prohibitively long input sequences involved~\citep{xue2022byt5}. Some prior studies have addressed this challenge by introducing more resource-efficient self-attention mechanisms~\citep{el2020characterbert, clark2022canine}, or by utilizing architectures that avoid attention entirely~\citep{wang2024mambabyte}~(\S\ref{section:related_work}). Nonetheless, these solutions are mainly effective for \textit{smaller models}. When scaling up, most of the computation in a Transformer is attributed to the extensive feed-forward network layers, which process every individual byte, rather than to the attention layers themselves.

To optimize compute allocation, we introduce a flexible, adaptive technique that organizes bytes into \textit{patches}~(\S\ref{section:patching}), alongside a novel architectural design that fuses information from both individual bytes and their aggregated patches. In contrast to tokenization methods, {\textsc BLT} does not rely on a predetermined patch vocabulary. Instead, any selection of byte groupings is transformed into latent patch embeddings through compact, trainable encoder and decoder components. Our findings demonstrate that this strategy enables a \textit{more} efficient use of computational resources compared to models utilizing tokenization.

Traditional tokenization-based large language models dedicate an equal amount of computational resources to every token, prioritizing uniformity at the expense of efficiency. This approach relies on compression techniques that may not align with the actual complexity of making predictions at different positions. Our proposed architecture instead emphasizes adaptive computation, directing resources to locations that require greater attention. For instance, accurately generating the final characters of a word is typically a straightforward, low-entropy prediction, whereas choosing the first token of a sentence tends to be much more challenging and information-rich, thereby demanding more computational investment. This principle underlies the design of {\textsc BLT}'s architecture~(\S\ref{section:architecture}), which incorporates three transformer modules: two compact, byte-level \textit{local models} and a larger, global \textit{latent transformer}~(\autoref{fig:arch_high_level}). 

To enable this selective allocation, {\textsc BLT} forms patches by clustering bytes according to the entropy associated with predicting subsequent bytes, resulting in context-aware groupings that maintain relatively constant information density.

We introduce the first flop-controlled scaling analysis of byte-level models trained end-to-end up to 8B parameters and 4T training bytes, achieving this without relying on fixed-vocabulary tokenization. In terms of training, {\textsc BLT} achieves parity with Llama 3’s flop-controlled performance\footnote{We measure a model’s compute expenditure as the total number of Floating Point OPerations (flops) performed.} while reducing inference flops by as much as 50\%~(\S\ref{section:scaling}). Additionally, direct processing of raw bytes leads to noticeable gains in modeling data distribution tails. Models using {\textsc BLT} show greater robustness against noisy data than those relying on tokenization and outperform them at capturing fine-grained information at the character level, with results validated on orthographic recognition, phonological reasoning, and machine translation for low-resource languages~(\S\ref{sec:robustness}). 

Moreover, {\textsc BLT} enables us to scale up both model and patch sizes together, all within a fixed inference flop constraint. Increasing patch size generally lowers computation per inference, freeing resources that can be invested to increase the global latent transformer’s capacity since it operates less frequently. Through inference-flop constrained scaling studies~(\autoref{fig:fixed_inference_scaling}), we find the scaling characteristics of {\textsc BLT} surpass those of models using token-based approaches.

In conclusion, this work offers the following main contributions: 1) We propose {\textsc BLT}, a byte latent LLM framework designed to allocate computational resources adaptively, resulting in enhanced FLOP efficiency. 2) We demonstrate that our method achieves training FLOP-controlled equivalence with Llama 3 up to the 8B parameter level, and further allows users to balance slight drops in evaluation scores against FLOP efficiency improvements as high as 50\%. 3) The {\textsc BLT} architecture enables a novel approach to scaling LLMs, permitting expansion in model size without exceeding a predetermined inference budget. 4) Our results highlight that {\textsc BLT} models show greater resilience to input noise and exhibit sensitivity to sub-word features in inputs, capturing information often overlooked by token-level LLMs. The source code for both training and inference with {\textsc BLT} is available at~\url{https://github.com/facebookresearch/blt}.

\section{Patching: From Individual Bytes to Groups of Bytes}
\label{section:patching}

Dividing bytes into \textit{patches} enables {\textsc BLT} to flexibly assign computational resources depending on the data's context. Figure~\ref{fig:patching_types} illustrates a range of strategies for splitting byte sequences into patches. More precisely, a patching function $f_p$ takes an input sequence of bytes $\pmb{x}=\{x_i,|i=1,\ldots n\}$, where the length is $n$, and partitions it into a sequence of $m < n$ patches $\pmb{p}=\{p_j|j=1,\ldots,m\}$ by designating each $x_i$ with a binary indicator, mapping to \{0,1\}; a value of 1 marks the initiation of a new patch. In both token-driven and patch-based architectures, the main Transformer’s computational expense scales with the number of main processing steps, which, in the case of {\textsc BLT}, equates to the patch count produced by a specific patching approach. Accordingly, the \textit{patch size}—the mean length of a patch—serves as the primary determinant for computational demand during training and inference given a selected patching scheme~(\S\ref{section:flops}).

We proceed to outline three types of patching functions: segmenting by a preset byte count per patch~(\S\ref{section:static-patch}), patching at whitespaces~(\S\ref{section:white-patch}), and patching adaptively using entropy signals from a lightweight byte-level language model~(\S\ref{section:dyn-patch}). We conclude by covering the concept of incremental patching and clarifying distinctions between patching and traditional tokenization~(\S\ref{section:bpe-patch}).

\subsection{Strided Patching Every K Bytes}
\label{section:static-patch}

One of the simplest approaches to grouping bytes involves dividing them into equally sized patches of length $k$, following the strategy adopted in MegaByte~\citep{yu2023megabyte}. This method, utilizing a constant stride, offers implementation simplicity during both training and inference phases, and provides an intuitive way to adjust the mean patch size, thereby offering straightforward control over computational requirements. Nevertheless, this strategy has notable limitations. It does not permit dynamic allocation of computational resources to regions where they are most critical: for instance, a transformer step $j$ may be inefficiently expended on predicting whitespace within code, or conversely, fail to devote adequate computation to segments of bytes containing dense informational content like mathematical expressions. Furthermore, this approach results in the inconsistent and context-agnostic segmentation of similar byte sequences, causing even identical words to be fragmented in different ways.

\subsection{Space Patching}
\label{section:white-patch}

\citet{slagle2024spacebyte} introduces an enhancement to strided patching by forming new patches immediately after encountering any space-like bytes\footnote{
    Space-like bytes are characterized as bytes which are neither latin characters, digits, nor utf-8 continuation bytes. Furthermore, at least one non space-like byte must be present in every patch.}, corresponding to common linguistic unit boundaries in many languages. The Space patching strategy dedicates a latent transformer step (implying increased computational cost) to each word, thus maintaining consistent patching of words across sequences and ensuring computational resources are focused on challenging predictions, which frequently occur right after spaces. For instance, generating the initial byte in the response to the question ``Who composed the Magic Flute?\uline{\hspace{.6em}}'' presents significantly higher complexity compared to generating the following bytes post-``M,'' since identifying the first character considerably narrows the field of plausible responses, rendering the rest of the completion ``Mozart'' substantially less ambiguous. Despite these strengths, Space patching struggles to flexibly accommodate all languages or domains and, crucially, does not allow for adaptable patch sizes. To address these issues, we propose a novel patching approach, motivated by the observation that initial bytes of words are generally most uncertain, and which introduces an effective and flexible patch size control mechanism.

\subsection{Entropy Patching: Using Next-Byte Entropies from a Small Byte LM}
\label{section:dyn-patch}

Instead of utilizing a rule-based heuristic like whitespace, our method adopts a data-driven strategy to detect points of high uncertainty in next-byte predictions. We propose \textit{entropy patching}, a technique that determines patch boundaries based on entropy estimates.

We train a small byte-level auto-regressive language model on the training data for {\textsc BLT} and compute next byte entropies under the LM distribution $p_{e}$ over the byte vocabulary $\mathcal{V}$:

\begin{equation}
H(x_i) = \sum_{v \in \mathcal{V}} p_{e}(x_i=v|\pmb{x}_{<i}) \log p_{e}(x_i=v|\pmb{x}_{<i})
\end{equation}

We explore two approaches for detecting patch boundaries using entropies $H(x_i)$. The first method selects points where entropy exceeds a fixed global threshold, as shown in~\autoref{fig:patching}. The second method looks for locations where the entropy is comparatively high when contrasted with the preceding value. This latter strategy can also be seen as finding points that disrupt a generally monotonically decreasing trend in entropy within a patch.

\begin{align*}
    \text{Global Constraint}\hspace{1cm}H(x_t)& > \theta_g\\
    \text{Approx. Monotonic Constraint}\hspace{1cm}H(x_t)& - H(x_{t-1}) > \theta_r
\end{align*}

Patch boundaries are determined in a streamlined preprocessing phase that takes place as data is being loaded. This contrasts with the approach of~\citet{nawrot-etal-2023-efficient}, who employ a classifier specifically trained to identify patch boundaries using entropy-based criteria. Within our experiments~(\S\ref{section:experiments}), we conduct a comparison between these two strategies for differentiating bytes of low and high entropy.

\subsection{The Byte-Pair Encoding (BPE) Tokenizer and Incremental Patching}
\label{section:incremental-patching}
\label{section:bpe-patch}

Numerous contemporary llms, such as our baseline Llama 3, employ a subword tokenization scheme like bpe~\citep{gage1994new, sennrich-etal-2016-neural}. In our context, ``tokens'' denote groups of bytes selected from a \textit{finite} vocabulary that is predefined before training, contrasting with ``patches,'' which indicate adaptively assembled sequences that do not rely on a preset vocabulary. An essential distinction is that, when using tokens, the model lacks direct exposure to the fundamental byte-level features.

An important advancement introduced by {\textsc BLT} relative to tokenization-based architectures is its new approach to balancing vocabulary size against computational requirements. In traditional language models, expanding the vocabulary typically results in larger average token lengths, thereby decreasing the number of decoding steps required. However, this also leads to an increase in the dimensionality of the model’s final projection layer, which raises the computational burden. Consequently, tokenization-based strategies have limited flexibility when it comes to varying token granularity or managing inference cost. To illustrate, Llama 3 achieves an increase in average token length from 3.7 to 4.4 bytes, yet this improvement requires scaling its embedding table by a factor of four compared to Llama 2.

During generation, {\textsc BLT} must determine at each point in the byte sequence whether it is at a patch boundary, as this choice impacts whether the Latent Transformer is employed for additional computation. Crucially, this boundary decision must be made without relying on any future, yet-to-be-generated segments of the sequence. Therefore, patching strategies cannot presume knowledge of subsequent bytes when partitioning the sequence. Mathematically, a patching function $f_p$ possesses the property of incremental patching if:
$$f_p(\pmb{x}_{<i}) = f_p(\pmb{x})_{<i}$$
bpe does not qualify as an incremental patching scheme, because it can segment identical prefixes differently depending on the following context, thus failing to meet the above criterion.\footnote{Introducing a dedicated delimiter token to signal patch boundaries can adapt bpe to function as an incremental patching scheme, but this comes at the cost of lengthening the byte sequence.}

\section{{\textsc BLT} Architecture}
\label{section:architecture}
The {\textsc BLT} system consists of a global, large-scale autoregressive language model that processes inputs in the form of patch representations. Complementing this, it utilizes a pair of lightweight local models: one responsible for converting byte sequences into patch representations, and the other for transforming these representations back into their original byte format~(\autoref{fig:arch_high_level}).

\subsection{Latent Global Transformer Model}

\textit{The Latent Global Transformer} refers to an autoregressive transformer architecture $\mathcal{G}$, composed of $l_{\mathcal{G}}$ layers, that transforms a sequence of input latent patch representations $p_j$ into corresponding output patch representations $o_j$. In this work, we employ the subscript $j$ to indicate patches, while $i$ refers to individual bytes. The global model implements a block-causal attention mask~\citep{dubey2024llama}, which confines the model’s attention span to all patches up to and including the present patch within each document. Notably, this transformer accounts for the majority of computational cost (in FLOPs) both during pre-training and inference. Consequently, the decision of when to utilize this model becomes a lever for allocating and modulating compute resources across different segments of the input and output, thereby adapting to varying levels of complexity.

\subsection{Local Encoder}
\label{section:local}

\textit{The Local Encoder Model}, represented as $\mathcal{E}$, utilizes a compact transformer architecture characterized by a number of layers $l_{\mathcal{E}}$ such that $l_{\mathcal{E}} << l_{\mathcal{G}}$. Its central purpose is to convert sequences of input bytes $b_i$ into informative patch representations, $p_j$, in an efficient manner. A significant innovation in this model, deviating from the standard transformer, is the integration of a cross-attention layer immediately after each transformer layer; this mechanism aggregates byte-level information into patch-level embeddings~(\autoref{fig:crossattn}).

The encoding process begins by embedding each byte $b_i$ with a matrix from $\mathbb{R}^{256 \times h_{\mathcal{E}}}$, resulting in $x_i$. These initial embeddings may be enhanced by concatenating hash-embeddings~(\S\ref{sec:hash-emb}), providing additional auxiliary information. Subsequently, the model applies a sequence of alternating transformer and cross-attention blocks~(\S\ref{sec:enc-cross-attn}), converting the initial byte-wise features into the patch representations, $p_i$, which are then forwarded to the global transformer $\mathcal{G}$.

The transformer stages employ a \textit{local block causal} attention strategy; this attention allows each byte to consider up to $w_{\mathcal{E}}$ prior bytes within a localized window. While these windows may encompass bytes straddling adaptive patch boundaries, they do not extend beyond the edges of individual documents. Details about the specific implementations of the embedding scheme and the cross-attention mechanism are provided in the subsequent sections.

\subsubsection{Encoder Hash n-gram Embeddings}
\label{sec:ngram-emb}
\label{sec:hash-emb}
To generate resilient and rich representations for each step $i$, it is essential to embed contextual information pertaining to previous bytes. {\textsc BLT} addresses this by encoding each byte $b_i$ not only in isolation but also within the context of a byte n-gram. At every position $i$, we initially generate byte-grams

\begin{align}
    g_{i,n}=\{b_{i-n + 1},\ldots, b_i\}
\end{align}

for each byte position $i$ and for $n$ ranging from three to eight.\footnote{
Byte-grams of length $n$ or more are excluded when $i < n$. }

Next, we propose hash $n$-gram embeddings, in which every byte $n$-gram is assigned—via a hashing operation—to a specific slot in an embedding table $E_{n}^{hash}$ with a predetermined capacity, for each $n$ in the set $\{3,4,5,6,7,8\}$~\citep{bai2010learning}. The corresponding embedding produced is added to the byte’s own embedding, after which it undergoes normalization and serves as the input to the local encoder. We compute the enhanced embedding as follows:

\begin{align}
e_i &= x_i + \sum_{n = 3,...,8}E_{n}^{hash}(\text{Hash}(g_{i,n})) \\
\text{where, Hash}(g_{i,n}) &= \text{RollPolyHash}(g_{i,n}) \% |E_{n}^{hash}|
\end{align}

We divide $e_i$ by the total count of $n$-gram sizes incremented by one, employing RollPolyHash as specified in Appendix~\ref{appendix:rollpolyhash}. Section~\ref{section:ablations} investigates how varying $n$-gram hash embedding parameters, including different choices of $n$ and embedding table dimensionality, influences scaling law patterns under a fixed computational budget. Alongside hashed $n$-gram embeddings, we explored embedding schemes based on $n$-gram frequency; further information regarding this approach is presented in Appendix~\ref{appendix:ngrams}.

\subsubsection{Encoder Multi-Headed Cross-Attention}
\label{sec:enc-cross-attn}
Our approach builds directly upon the cross-attention mechanism utilized in the input module of the Perceiver architecture~\citep{jaegle2021perceiver}, with a key distinction: in our method, the latent representations are associated with individual patch representations that can vary, rather than relying on a fixed collection of latents~(\autoref{fig:crossattn}), and each latent is restricted to attend solely to the bytes forming its specific patch. In this setup, the module generates a query vector for each patch $p_j$, where initialization is accomplished by aggregating the byte-level representations within patch $p_j$—for example, via pooling—subsequently processed through a linear projection layer, $\mathcal{E}_{C} \in \mathbb{R}^{h_{\mathcal{E}} \times (h_{\mathcal{E}} \times U_{\mathcal{E}})}$, with $U_{\mathcal{E}}$ representing the encoder cross-attention head count. Specifically, let $f_{\text{bytes}}(p_j)$ refer to the sequence of bytes for patch $p_j$; the following computation ensues

\begin{align}
P_{0,j} &= \mathcal{E}_{C}(f_\text{bytes}((p_j)), f~ \text{is a pooling function} \\
P_l &= P_{l-1} + W_o\left(\text{softmax}\left(\frac{QK^T}{\sqrt{d_k}}\right)V\right) \\
\text{where } Q_j &= W_q(P_{l-1,j}), K_i = W_k(h_{l-1,i}), V_i = W_v(h_{l-1,i}) \\
h_l &= \text{Encoder-Transformer-Layer}_l(h_{l-1})
\end{align}

Here, $P \in \mathbb{R}^{n_p \times h_{\mathcal{G}}}$ denotes the $n_p$ patch representations that will serve as inputs to the global model. These are constructed by aggregating the byte embeddings $e_i$ associated with each patch $p_j$ through a pooling operation. The matrices $W_q$, $W_k$, $W_v$, and $W_o$ correspond to the linear projections used to compute queries, keys, values, and outputs, respectively. In this setup, the keys and values are computed as projections of the byte representations $h_i$ output by the preceding layer (or $e_i$ when processing the initial layer). To restrict information flow within patches, a specialized masking mechanism is applied such that each query $Q_j$ only attends to those keys and values originating from the bytes present in patch $j$. Since multi-head attention is employed over $Q$, $K$, and $V$, and given that the dimensionality of the patch representations ($h_{\mathcal{G}}$) usually exceeds $h_{\mathcal{E}}$, we keep $P_l$ in the form of several heads, each with dimension $h_{\mathcal{E}}$, during the cross-attention step, and subsequently concatenate these into the $h_{\mathcal{G}}$-dimensional representation. In this architecture, a pre-LayerNorm normalization is applied to queries, keys, and values, and positional embeddings are omitted from the cross-attention layer. The overall design incorporates a residual connection wrapping the cross-attention module.

\subsection{Local Decoder} 

The local decoder $\mathcal{D}$, analogous to the local encoder, utilizes a compact transformer architecture composed of $l_{\mathcal{D}}$ layers, where $l_{\mathcal{D}} << l_{\mathcal{G}}$. Its primary task is to reconstruct the original raw byte sequence $y_i$ from the global patch representations $o_j$. The prediction of each byte within the sequence depends on the context provided by preceding bytes, leveraging the hidden states output by the local encoder corresponding to the input byte sequence. The local decoder architecture interleaves $l_{\mathcal{D}}$ cross-attention layers and transformer blocks. Specifically, each cross-attention mechanism in the decoder precedes the transformer layer, mapping patch representations to byte-level representations, after which the local decoder's transformer block further processes the byte sequence.

\subsubsection{Decoder Multi-headed Cross-Attention} 

In the decoder cross-attention, the roles of the queries and key/values are interchanged i.e. the byte-representations are now the queries, and the patch representations are now the key/values. The initial byte-representations for the cross-attention are initialized as the byte embeddings from the last encoder layer i.e. $h_{l_{\mathcal{E}}}$. The subsequent byte-representations for layer $l$, $d_{l,i}$ are computed as:

\begin{align}
D_0 &= h_{l_{\mathcal{E}}} \\
B_l &= D_{l-1} + W_o\left(\text{softmax}\left(\frac{QK^T}{\sqrt{d_k}}\right)V\right), \\
  &\text{where } Q_i = W_q(d_{l-1,i}), K_i = W_k(\mathcal{D}_C(o_j)), V_i = W_v(\mathcal{D}_C(o_j)) \\
  D_l &= \text{Decoder-Transformer-layer}_l(B_l)
\end{align}

In this setup, $W_k$ and $W_v$ denote the key and value projection matrices, respectively. These matrices act on the output patch representations $o_j$ produced by the global model, via a linear transformation followed by a split operation denoted by $\mathcal{D}_C$. Meanwhile, $W_q$ refers to the query projection matrix applied to byte-level representations $d_{l-1}$ from the preceding layer of the decoder transformer, or to $h_{l_{\mathcal{E}}}$ if it is the initial layer. The matrix $W_o$ serves as the output projection, resulting in $B$ having the shape $\mathbb{R}^{h_{\mathcal{D}} \times n_b}$, where $n_b$ corresponds to the total number of predicted output bytes. The subsequent decoder outputs $D_l$ are generated by applying a decoder transformer layer to the result of the cross-attention module, $B$. The cross-attention block—analogous to that used in the local encoder—utilizes multi-head attention, pre-LayerNorm normalization, omits positional encodings, and incorporates a residual pathway around the module.

\section{Experimental Setup}
\label{section:experiments}
Our experimental design rigorously controls for confounding factors in order to equitably evaluate {\textsc BLT} alongside models that rely on tokenization. Special care is taken to prevent {\textsc BLT} from benefitting from any inadvertent increase in sequence length or context relative to the baseline tokenization-based systems.

\subsection{Pre-training Datasets}

In this study, all model sizes are initially pre-trained using two principal datasets: (1) the Llama 2 dataset~\citep{touvron2023llama}, encompassing 2 trillion tokens sourced from diverse, publicly accessible materials, which are rigorously cleaned and filtered to ensure higher data quality; and (2) {\textsc BLT}-1T, a new compilation consisting of 1 trillion tokens obtained from multiple public sources, including a selected fraction of the pre-training set made available by Datacomp-LM~\citep{li2024datacomp}. The Llama 2 corpus is leveraged for scaling law analyses, especially focusing on the optimal token count as identified by~\cite{dubey2024llama}, informing the architectural design of {\textsc BLT}. In contrast, {\textsc BLT}-1T serves as the main corpus for end-to-end pre-training in order to benchmark against Llama 3 performance on downstream evaluation tasks. Notably, both corpora are devoid of any data derived from Meta platforms or services. Additionally, when conducting baseline comparisons for models relying on tokenizers, we employ the Llama 3 tokenizer, which features a 128K token vocabulary; empirical findings from our work indicate it consistently delivers superior baseline results compared to the Llama 2 tokenizer.

\subsection{Entropy Model} In our study, the entropy model employed is a byte-level language model, and it is trained on the identical distribution as the complete {\textsc BLT} model. By default, unless specified otherwise, we utilize a transformer architecture comprising 100M parameters distributed across 14 layers, each with a hidden size of 512. The model uses sliding window attention spanning 512 bytes. All other hyperparameters are aligned with those used in the local and global transformer baselines. Various configurations—including different model capacities, receptive fields, and model architectures—were evaluated, as elaborated in \autoref{section:ablations}. Notably, for models with sufficiently limited receptive fields, the learned entropy model parameters can be efficiently represented using a lookup table.

\subsection{Entropy Threshold and Equalizing Context Length} For models implementing entropy-based patching, we determine a patching threshold to yield a specified average \textit{patch size} across the pretraining data mix. In {\textsc BLT}, as opposed to conventional tokenization methods, the \textit{patch size} is selectable without restrictions, which substantially impacts the model’s context length. To prevent models with larger patch sizes from receiving undue benefit due to increased context, we control for the expected total number of bytes per batch, keeping it fixed. Consequently, sequence lengths for models with larger patches are proportionally shortened. Specifically, on Llama 2 data, we apply an 8k byte context window, while for the {\textsc BLT}-1T dataset, we expand this to a 16k byte context on average, ensuring a constant average batch size of 16M bytes in both cases.

Although the mean batch size remains fixed, dynamic patching approaches result in varying proportions of bytes per patch when assembling data batches. In our implementation, {\textsc BLT} training composes batches based on patches to eliminate the need for padding within the computationally intensive latent transformer, thereby maintaining a uniform patch count across all batches. To prevent memory usage surges caused by exceptionally large patches, byte sequences are padded or, if necessary, truncated to lengths of 12k and 24k bytes for Llama 2 and {\textsc BLT}-1T datasets, respectively, during the training process.

\subsection{Entropy Model Context}
\label{section:entropy-context}
Our experiments indicate that entropy patching leads to increasingly larger patches in highly structured content, such as in multiple choice questions (refer to the MMLU example in \autoref{fig:mmlu_patching}). Such formats tend to be highly repetitive. The observed variance is attributable to the reduced entropy present in the repeated segments within the entropy model context. Consequently, for the large-scale {\textsc BLT}-Entropy experiment using a patch size of 4.5, we refresh the entropy context by inserting new lines and implement an approximate monotonicity constraint, which demonstrates greater robustness to "entropy drift" as context length fluctuates. It should be noted that this adjustment modifies only the entropy calculation process; our method for determining the entropy threshold remains unchanged.

\subsection{FLOPs Estimation} 
\label{section:flops}

Our methodology for calculating transformer FLOPs is primarily based on the formulations described in Chinchilla~\citep{hoffmann2022training}, which account for the FLOPs associated with the feed-forward modules, qkvo projections in self-attention mechanisms, attention computation, and the final output projection. One key distinction in our approach is the assumption that the input embedding layer operates via a computationally efficient lookup rather than through dense matrix multiplication, resulting in a negligible, effectively zero, FLOP cost for this component. Consistent with earlier research, we assume that the backward pass requires twice as many FLOPs as the forward pass.

To compute flops \textit{per byte} for {\textsc BLT} models, we add up the flops for the local encoder transformer, the global latent transformer, and the local decoder transformer, together with the cross attention blocks in the encoder and the decoder:

\begin{align}
    \text{FL}_{\text{{\textsc BLT}}} &= \text{Transf. FL}(h_{\mathcal{G}}, l_{\mathcal{G}}, m=n_{ctx}/n_p,V=0) / n_p\\
   &+ \text{Transf. FL}(h_{\mathcal{E}}, l_{\mathcal{E}}, m=w_{\mathcal{E}}, V=0) \\
   &+ \text{Transf. FL}(h_{\mathcal{D}}, l_{\mathcal{D}}, m=w_{\mathcal{D}}, V=256) \\ 
    &+ \text{Cross Attn. FL}(h_{\mathcal{E}}, l_{\mathcal{E}}, m=n_p, r=n_p/k) \cdot \frac{k}{n_p} \\ 
   &+ \text{Cross Attn. FL}(h_{\mathcal{D}}, l_{\mathcal{D}}, m=k, r=k/n_p)
\end{align}

In this context, $n_{ctx}$ denotes the input sequence length measured in bytes, while $n_p$ refers to the size of the patches. The parameter $r$ specifies the proportion of queries relative to keys and values, and $k$ captures the relationship between the patch dimension and the byte dimension—in other words, it is the count of local model segments concatenated to establish a unified model representation ($k=2$ in \autoref{fig:crossattn}). The variable $V$ represents the output projection’s vocabulary size, a parameter exclusive to the local decoder. The context length for attention, denoted as $m$, varies depending on whether the module processes byte-level or patch-level sequences. We adjust our calculations of attention FLOPs for each architectural component accordingly. Detailed formulations for the computation of FLOPs in both Transformer and Cross-Attention settings can be found in Appendix~\ref{section:flops-appendix}.

\subsection{Bits-Per-Byte Estimation} Perplexity only makes sense in the context of a fixed tokenizer as it is a measure of the uncertainty for each token. When comparing byte and token-level models, following previous work~\citep{xue2022byt5, yu2023megabyte, wang2024mambabyte}, we instead report Bits-Per-Byte (BPB), a tokenizer independent version of perplexity. Specifically:

\begin{align}
    \text{BPB}(x)&= \frac{\mathcal{L}_{CE}(\pmb{x})}{\ln(2)\cdot n_{\text{bytes}}}
\end{align}

Here, the total cross-entropy loss for the data $\pmb{x}$, which quantifies the model's uncertainty, is divided by both the number of bytes in $\pmb{x}$ and a normalization constant.

\subsection{Transformer Architecture Hyperparameters} In configuring the transformer blocks for both local and global models within {\textsc BLT}, we primarily adhere to the Llama~3~\citep{dubey2024llama} design. The feed-forward components incorporate the SwiGLU activation function~\citep{swiglu}, and rotary positional embeddings (RoPE)~\citep{RoPe} with $\theta=500000$~\citep{xiong2024effective} are applied exclusively to self-attention mechanisms. For normalization, RMSNorm~\citep{RMSNorm} is employed across all layers. Self-attention operations utilizing fixed-standard attention masks—such as \textit{block causal} and \textit{fixed-window block causal}—leverage Flash attention~\citep{dao2022flashattention}, where a window size of 512 is set for fixed-width masking. For cross-attention layers, which require dynamic and patch-specific attention masks, we deploy Flex Attention\footnote{\url{https://pytorch.org/blog/flexattention}} to take advantage of optimized fused kernels, markedly enhancing training efficiency.

\subsection{{\textsc BLT}-Specific Hyperparameters} 

Our investigation into the performance of {\textsc BLT} models comprises experiments centered on scaling behavior and evaluation on downstream tasks, examining model sizes of 400M, 1B, 2B, 4B, and 8B. Appendix~Table~\ref{tab:arch} details the corresponding architectural hyperparameters. For the initial cross-attention layer within the local encoder, queries are initialized via max-pooling. Across all {\textsc BLT} configurations, we utilize $500,000$ hashes with a single hash function and n-gram windows spanning 3 to 8. A uniform learning rate of $4e-4$ is maintained for every model instance. The equivalence of optimal learning rates between token models and {\textsc BLT} models is established through a search across $1e-3$ to $1e-4$ at the 400M and 1B scales, indicating that both benefit from the same rate. For scaling experiments on Llama-2 data, the batch sizes employed adhere to the recommendations from~\cite{dubey2024llama}, matched by data volume when necessary. For the optimization procedure, the AdamW optimizer~\citep{adamw} is employed with $\beta_1 = 0.9$, $\beta_2 = 0.95$, and $\epsilon = 10^{-8}$. The learning rate schedule comprises a linear warm-up phase over 2000 steps, followed by cosine decay down to zero. Regularization employs a weight decay of 0.1 and global gradient clipping with a cap of 1.0.

\section{Scaling Trends}
\label{section:scaling}

We provide a comprehensive overview of how byte-level models scale, with the goal of guiding further advancements of {\textsc BLT} models. Our exploration of scaling is designed to overcome certain shortcomings in earlier studies on byte-level approaches in these key ways: (a) We systematically analyze performance within a compute-optimal training context, (b) We construct equivalent 8B parameter models utilizing substantial volumes of data—reaching up to 1T tokens (4T bytes)—and assess their capabilities on subsequent downstream evaluations, and (c) We scrutinize scaling properties under settings where inference cost is kept constant. Later, we examine the particular strengths that arise from processing byte sequences directly.

\subsection{Parameter Matched Compute Optimal Scaling Trends}
\label{section:parameter-matched-scaling}
Leveraging the Llama 2 dataset, we conduct experiments by training a collection of \textit{compute-optimal} bpe and {\textsc BLT} models at four distinct scales, with parameter counts ranging from 1B to 8B. We assess their training compute, measured in flops, in relation to language modeling outcomes using a representative fraction of the training data mixture. For the bpe models, we follow the optimal model size to data ratio as prescribed by Llama~3~\citep{dubey2024llama}, ensuring alignment with \textit{compute-optimal} training objectives designed to maximize performance on the available dataset given a fixed computational budget~\citep{hoffmann2022training}. This approach allows these bpe models to serve as a rigorous performance reference. Each bpe model is paired with a {\textsc BLT} counterpart, trained on the identical data split and constructed using a Latent Transformer that mirrors the corresponding bpe Transformer's size and architecture.

As shown in~\autoref{fig:scaling-compute-opt} (right), {\textsc BLT} models consistently achieve performance that is comparable to or better than that of their bpe equivalents, and this relationship persists with increases in both model size and computational cost. To our knowledge, {\textsc BLT} represents the first byte-level Transformer framework to reach parity in scaling behaviors with BPE-based architectures under compute-optimal conditions. These findings support our hypothesis that the optimal parameters-to-training compute ratio established for bpe models is also applicable to {\textsc BLT}, or at least does not differ significantly.

Both advancements in architecture and the use of dynamic patching are essential to align with bpe scaling patterns. In ~\autoref{fig:scaling-compute-opt} (left), we present a comparison between models utilizing space-patching techniques and Llama 3. To approximate SpaceByte \citep{slagle2024spacebyte}, we implement {\textsc BLT} space-patching but omit n-gram embeddings and cross-attention. While SpaceByte demonstrates gains over Megabyte, a significant performance gap persists compared to Llama 3. The impact of architectural enhancements in conjunction with dynamic patching is depicted in \autoref{fig:scaling-compute-opt} (right). At scale, {\textsc BLT} models achieve results comparable to leading tokenizer-based architectures like Llama 3.

Additionally, our results indicate that the selection of tokenizer has a tangible impact on the performance of tokenizer-based models; specifically, models utilizing the Llama-3 tokenizer demonstrate superior results compared to those employing the Llama-2 tokenizer, when both are trained on identical datasets.

In summary, our {\textsc BLT} architecture demonstrates a scaling pattern that falls between that of Llama 2 and Llama 3 when employing much larger patch sizes. The average token lengths for the bpe tokenizers in Llama 2 and 3 are 3.7 and 4.4 bytes, respectively. However, {\textsc BLT} is able to replicate comparable scaling behaviors when the average patch size is extended to 6 or even 8 bytes. Because inference flops decrease as the average patch size increases, adopting a patch size of 8 bytes results in an approximate 50\% reduction in inference flops. Furthermore, models that use larger patch sizes tend to exhibit improved performance as both the model and dataset scale up. For instance, while {\textsc BLT} with an 8-byte patch size initially underperforms relative to bpe Llama 2 at the 1B parameter mark, it surpasses bpe by the 7B scale. This pattern implies that models could benefit even more from larger patch sizes at greater scales, and that employing larger patches may become increasingly viable as model capacity and available computational resources expand.

\subsection{Beyond Compute Optimal Task Evaluations}
\label{section:evals}
To further investigate scaling behaviors, we train an 8B {\textsc BLT} model exceeding the compute-optimal scale on the {\textsc BLT}-1T dataset—a larger and higher-quality corpus—and evaluate its capabilities across an array of widely-used classification and generative tasks. The assessment encompasses benchmarks that test common sense reasoning, general world knowledge, and code generation proficiency:
\paragraph{Classification tasks} consist of ARC-Easy (0-shot)~\citep{Eval_ARC}, Arc-Challenge (0-shot)~\citep{Eval_ARC}, HellaSwag (0-shot)~\citep{Eval_hellaswag}, PIQA (0-shot)~\citep{Eval_piqa}, and MMLU (5-shot)~\citep{hendrycks2020measuring}. To evaluate these, we utilize a prompt-based likelihood scoring approach on candidate answers and present mean accuracy across the tasks.
\paragraph{Coding related generation tasks:} We assess model performance with pass@1 on MBPP (3-shot)~\citep{Eval_MBPP} and HumanEval (0-shot)~\citep{Eval_HumanEval}, measuring each model's competence in generating executable Python code.

In \autoref{tab:evals}, we present a comparison among three models, all trained on the {\textsc BLT}-1T dataset: a bpe model that utilizes the Llama 3 tokenizer\footnote{The Llama 3 tokenizer, featuring a 128k vocabulary, is selected due to its superior performance relative to the Llama 2 tokenizer, which has a 32k vocabulary.}, and two different versions of the {\textsc BLT} architecture. Specifically, one variant applies a space-based patching method ({\textsc BLT}-Space), while the other adopts an entropy-driven patching approach ({\textsc BLT}-Entropy), in which approximate monotonicity is enforced and the model context is reset at each newline character, as detailed in~\autoref{section:entropy-context}. Each model is trained with a comparable computational (flop) budget. Notably, for {\textsc BLT}-Entropy, we further adjust the entropy threshold at inference from 0.6 to 0.1; this modification leads to improved performance on evaluation tasks, albeit requiring a greater number of inference steps.

The {\textsc BLT}-Entropy model achieves superior results compared to the Llama 3 model on four out of seven benchmarks, despite both being trained with an equivalent data volume in terms of bytes. This performance gain can likely be attributed to two factors: (1) more efficient utilization of training compute enabled by dynamic patching, and (2) explicit modeling at the byte level rather than relying on token representations.

Conversely, {\textsc BLT}-Space demonstrates inferior performance relative to the Llama 3 tokenizer across nearly all tasks, except for one. Nevertheless, due to its increased average patch size of 6 bytes, it results in a notable decrease in inference flops. For comparison, both the bpe and entropy-patching models exhibit an average patch size of about 4.5 bytes over the mixed training dataset. When provided with an equivalent training budget, the model employing larger patch sizes is able to process 30\% more data than the other two approaches. This expanded coverage, however, may drive {\textsc BLT} farther from the compute-optimal regime.

\subsection{Patches Scale Better Than Tokens} The {\textsc BLT} models offer the flexibility to scale both the architecture and the patch size in tandem, all while preserving a fixed computational cost for training and inference and keeping the training dataset unchanged. This ability to arbitrarily enlarge patch sizes is particular to patch-based architectures, allowing them to circumvent the efficiency limitations imposed by the fixed-vocabulary of token-based methods, as outlined in Section~\ref{section:bpe-patch}. Expanding patch sizes reduces computational requirements, thereby making resources available to increase the capacity of the global latent transformer, since this component is invoked less frequently.

We carry out a fixed inference scaling investigation to examine whether utilizing larger models that require fewer inference steps on increased patch sizes can surpass the performance of smaller models necessitating more steps. Beginning with baseline models containing 400 million and 3.6 billion parameters utilizing the Llama 2 tokenizer, we identify flop-equivalent configurations for both the Llama 3 tokenizer and {\textsc BLT}-Entropy models, assigning average patch sizes of 6 and 8 bytes across the training datamix (refer to~\autoref{tab:fixed-inf-params} for comprehensive model specifications). For those models with a patch size of 8, we opt to use 3 encoder layers instead of just 1. Each model is trained under several different training flop constraints.

\autoref{fig:fixed_inference_scaling} illustrates that {\textsc BLT} models display superior scaling behavior compared to tokenization-based models across both evaluated inference flop categories. While BPE architectures tend to yield higher performance at lower training compute expenditures, {\textsc BLT} models soon overtake them—shortly past the compute-optimal point. This suggests that, from a practical standpoint, allocating greater resources to initial pretraining can result in a model that delivers stronger performance under a set inference compute constraint. The family of 8B models, such as Llama 3.1, exemplifies this approach, as they have been trained on datasets nearly two orders of magnitude larger than those dictated by the compute-optimal guideline for their parameter count.

The threshold at which {\textsc BLT} surpasses token-based approaches has moved nearer to the compute-optimal regime when considering models with greater flop capacity, specifically shifting from being 3 times to 2.5 times above the compute-optimal limit. Furthermore, the model with a larger patch size of 8 displays a more pronounced scaling trajectory in this higher flop regime, surpassing alternative models earlier. As elaborated in Section~\ref{section:parameter-matched-scaling}, using larger patch sizes tends to yield performance that more closely approaches BPE models at increased model sizes. We partly explain this observation by the fact that the proportion of total flops attributed to the byte-level Encoder and Decoder components—whose scaling appears slower relative to the Latent Transformer—diminishes as models increase in size. When increasing total parameters twentyfold, from 400M to 8B, the local parameters specific to {\textsc BLT} expand by only about two times. This distinction is noteworthy since increasing patch size impacts only the flops in the patch Latent Transformer rather than those in the byte-level modules. Accordingly, the model size of {\textsc BLT}-Entropy with patch size 8 moved from 1.6 times to 1.7 times that of Llama 2 upon scaling to larger model configurations.

To conclude, our investigation into patch-length scaling reveals that the {\textsc BLT} patch-based framework exhibits superior scaling behavior when both patch size and model capacity are jointly enlarged. These favorable scaling patterns not only continue but may even become more pronounced with further increases in model scale.

\section{Byte Modeling Improves Robustness}
We further evaluate the robustness of {\textsc BLT} relative to token-based counterparts that do not utilize byte-level representations, and we introduce a method to augment pretrained token-based models with byte-level capabilities.
\label{sec:robustness}

\subsection{Character-Level Tasks}

One of the initial incentives for developing models that operate at the byte level was their enhanced resilience to noise present in the input at the byte scale, as well as their capability to recognize the subcomponents within tokens—a task where traditional tokenizer-based models often encounter challenges. To assess these characteristics, we conduct supplementary evaluations using benchmarks designed to test both the models' robustness against noise in the input and their sensitivity to characters. These benchmarks encompass English and multilingual characters, digits, and phonemes. The outcomes of these assessments are reported in Table~\ref{tab:char_tasks}.

\paragraph{Noisy Data} To assess robustness, we generate altered variants of the standard classification datasets referenced in Section~\ref{section:evals}, enabling a direct comparison between tokenizer-based models and {\textsc BLT}. Five separate character-level perturbation methods are used to induce diversity in the input: (a) \textit{AntSpeak}: Reformats the text so that every character is uppercase and separated by spaces. (b) \textit{Drop}: Randomly deletes 10\% of characters from the original text. (c) \textit{RandomCase}: Randomizes letter casing by assigning half of the characters to uppercase and the other half to lowercase. (d) \textit{Repeat}: Selects 20\% of characters to be repeated, up to four repetitions each. (e) \textit{UpperCase}: Converts all alphabetic characters to uppercase. In our evaluations, each perturbation is independently applied to the prompt, the completion, or to both, with performance results averaged over these variations. As shown in Table \ref{tab:char_tasks}, on the noised HellaSwag~\citep{Eval_hellaswag} dataset, {\textsc BLT} consistently demonstrates greater resilience than tokenizer-dependent models, exhibiting an average lead of 8 points compared to models trained on equivalent data, and even surpassing the Llama 3.1 model that was trained on a far larger corpus.

\paragraph{Phonology - Grapheme-to-Phoneme (G2P)} We evaluate the proficiency of {\textsc BLT} in converting grapheme sequences (letters comprising a word) into their corresponding phoneme transcriptions (which reflect pronunciation). Table \ref{tab:char_tasks} summarizes the 5-shot G2P performance using the Phonology Bench~\citep{suvarna2024phonologybench}. Our findings indicate that {\textsc BLT} achieves superior results compared to the baseline model relying on the Llama 3 1T tokenizer for this task.

\paragraph{CUTE}
To measure character-level comprehension, we test {\textsc BLT} on the CUTE benchmark~\citep{edman2024cute}, which consists of various tasks divided into three general categories: compositional understanding, orthographic similarity, and sequence manipulation. CUTE is particularly demanding for models that rely on tokenization, as these models often know the internal spelling of their tokens but find it difficult to leverage that knowledge in actual text manipulation. As shown in Table~\ref{tab:char_tasks}, {\textsc BLT}-Entropy achieves results more than 25 points higher than the BPE Llama 3 variants on CUTE. The model especially excels in character-level manipulation, reaching 99.9\% accuracy on both spelling-related tasks. These substantial gains are particularly noteworthy considering {\textsc BLT} was trained with only one-sixteenth the data used by Llama 3.1, suggesting that BPE-based models find it inherently challenging to acquire fine-grained character-level knowledge. Figure \ref{fig:cute_examples} highlights examples in which the Llama 3 tokenizer-based approach underperforms while our {\textsc BLT} demonstrates strong performance. The only areas where BPE-based methods retain an advantage are in word-level deletion and insertion; while these word-centric manipulations may not be as intuitively accessible to a byte-oriented model, the performance gap remains moderate, and constructing word representations from character-level information may prove less difficult than decomposing word units into characters. All experiments utilize an identical evaluation protocol and incorporate the original Huggingface prompts. Notably, BPE approaches could potentially be improved through specialized prompt engineering.

\paragraph{Low Resource Machine Translation} We assess the performance of {\textsc BLT} on translation tasks involving twenty-one low-resource languages spanning six widely represented language families and multiple scripts, leveraging the FLORES-101 benchmark~\citep{goyal2022flores}. SentencePiece BLEU scores are detailed in Table~\ref{tab:flores}. Our findings show that {\textsc BLT} achieves superior results compared to a baseline model using the Llama 3 tokenizer, with a 2-point overall improvement for English target translations and a 0.5-point gain for English source translations. For widely spoken language pairs, the results of {\textsc BLT} are on par with or marginally exceed those obtained by Llama 3. In contrast, for many language pairs within under-resourced language families, {\textsc BLT} delivers stronger outcomes than Llama 3, highlighting the capacity of byte-level modeling to effectively generalize to diverse and rare byte sequences.

\subsection{Training {\textsc BLT} from Llama 3}
\label{sec:distillation}
We investigate a method where {\textsc BLT} models benefit from already established tokenizer-based pre-trained models to improve both the speed and quality of training convergence. Specifically, this is done by initializing {\textsc BLT}'s global transformer parameters with those from a Llama 3.1 checkpoint. After initialization, we fine-tune the global transformer parameters using a learning rate that is one-tenth of that utilized for both the local encoder and local decoder, following the same number of training steps that is optimal for Llama 3. The resulting performance is then compared to a standard {\textsc BLT} baseline in Table~\ref{tab:distill}. Results clearly demonstrate that the {\textsc BLT} model initialized from Llama 3.1 consistently outperforms both the Llama 3 and original {\textsc BLT} baselines under equivalent computational budgets. Notably, even when contrasted with our {\textsc BLT}-Entropy variant—which was trained on a dataset approximately an order of magnitude larger (1T tokens)—the {\textsc BLT} initialized from Llama 3.1 achieves even higher scores on the MMLU benchmark. This finding indicates the potential of this initialization strategy to considerably reduce the computational requirements for training.

This approach can be interpreted as converting models that rely on tokenizers into models that do not, in essence transforming a pre-trained LLaMA 3.1 model into a {\textsc BLT} architecture. For a thorough comparison, we report results for the original LLaMA 3.1 model, which was trained on 15T tokens, in Table \ref{tab:distill}, and assess its performance relative to the {\textsc BLT} version adapted from LLaMA 3. We observe that while there is only a minor reduction in performance on MMLU and HumanEval, the degradation is more substantial for other benchmarks. This indicates that additional investigation is necessary to maximize the benefits of the pre-trained model, especially regarding the refinement of training data combinations and adjustment of hyperparameters.

\section{Ablations and Discussion}
\label{section:ablations}
Here, we present a set of ablation studies that evaluate and rationalize our design decisions for the {\textsc BLT} architecture, along with our selection of patching methods and hyper-parameter configurations used in training the {\textsc BLT} model with 8B parameters on the {\textsc BLT}-1T dataset.

\paragraph{Entropy Model Hyper-parameters} To investigate how entropy model size and context window length impact scaling behavior, we experiment with byte-level entropy transformer models ranging from 1m to 100m parameters, adjusting the context window length between 64 and 512. We display the scaling laws of bpb versus training flop, utilizing our $400m$ and $1b$ {\textsc BLT} models that were trained on the Llama-2 dataset, as shown in \autoref{fig:entropy_ablation}. The results indicate that performance improvements generally track increases in both entropy model size and context window length, although gains taper off once the model exceeds 50m parameters.

\paragraph{Types of Patching}

We evaluate the impact of the four patching methods described in Section~\ref{section:patching}, namely: 1) Strided Patching with strides of 4 and 6, 2) Whitespace-based patching, 3) Patching guided by the BPE Tokenizer utilizing the Llama 3 tokenizer, and 4) Entropy-driven patching implemented with a compact byte-level language model.

Although dynamic patching shortens the apparent sequence length, we ensure that context length remains comparable across different patching methods by holding the sequence length constant. Consequently, across training and inference, every model is exposed to an equivalent number of bytes per sequence on average, thereby ruling out the influence of increased context modeling as a confound. Figure~\ref{fig:scaling-compute-opt} presents the findings from these ablation studies. Each of the alternative patching strategies surpasses static patching, with space patching showing nearly equivalent performance to dynamic entropy-based patching.

In \autoref{tab:evals_ablation}, we provide evaluation results for the {\textsc BLT} models, highlighting performance comparisons among tokenizer-based approaches, space patching, and entropy-based patching. All models were trained on the Llama 2 dataset, using an optimal training duration as detailed in~\citep{dubey2024llama}. While space patching employs a more straightforward methodology that does not require the deployment of an entropy model dynamically during training, our results indicate that the benefits seen from entropy-based patching in scaling experiments~(Section~\ref{section:scaling}) are also present in broader downstream evaluation tasks.\footnote{The reported space patching outcomes were obtained from initial runs without cross-attention; however, analogous patterns emerge when cross-attention is utilized.}

\paragraph{Cross-Attention} 
In~\autoref{tab:crossattn_ablations_1b}, we present an ablation study examining the impact of cross-attention mechanisms incorporated at different stages within both the encoder and decoder components of {\textsc BLT}. Regarding encoder cross-attention, we consider several initialization strategies for the queries: (1) employing an identical learned embedding across all global states, (2) utilizing a hash embedding derived from the bytes of the patch, and (3) applying pooling to the encoder's hidden representation of patch bytes at the respective encoder layer.

Our results indicate that implementing cross-attention within the \textit{decoder} yields the greatest benefit. Introducing cross-attention in the encoder offers a modest gain, which manifests primarily when the queries are initialized via pooling. Furthermore, cross-attention proves advantageous especially on the Common-Crawl dataset, with its effectiveness becoming more pronounced as patch sizes increase.

\paragraph{n-gram Hash Embeddings} 

We experiment with n-gram hash embedding vocabularies of size 0, 100K, 200K, and 400K, with the corresponding outcomes summarized in Table~\ref{tab:ngram_ablations_1b}. Our analysis demonstrates that incorporating hash embeddings yields improvements across all datasets, most notably for Wikipedia and Github, where a 0.04 bpb gain is observed, compared to only a 0.01 bpb increase after 15k training steps at the 8B parameter scale. When adjusting the number of hashes from 500K down to 300K at 8B, the impact on performance after 15k steps is minimal, at approximately 0.001 bpb. These findings suggest that hash embeddings play a crucial role in elevating the performance of {\textsc BLT} to be comparable with tokenizer-based counterparts; nonetheless, performance benefits plateau beyond 300K hashes. Moreover, enhancements from hash embeddings appear to be largely orthogonal to those gained from cross-attention, as each technique predominantly benefits different datasets.

\paragraph{Local Model Hyperparamaters} Table \ref{tab:local_layers} presents an ablation study exploring different configurations for the number of layers in both the local encoder and decoder. The results indicate that when utilizing hash n-gram embeddings, {\textsc BLT} achieves strong performance even with a highly simplified encoder—specifically, one comprising a single layer—while the decoder benefits from having additional layers.

\section{Related Work}
\label{section:related_work}

\textbf{Character-Level RNNs:} The task of Character Language Modeling has held continued interest since the advent of early neural architectures~\citep{sutskever2011generating,mikolov2012subword,graves2013generating}, largely due to its inherent ability to handle out-of-vocabulary words natively, avoiding the need for explicit back-off strategies. In their work, \cite{kim2016character} develop a system that exclusively ingests characters on the input side by leveraging convolutional networks followed by highway layers before passing features to an LSTM-based RNN. This design achieved parity with the contemporary state-of-the-art RNN language models for English and demonstrated superior results for languages with complex morphology, underscoring another notable benefit of character-level LLMs. Expanding upon this, \cite{kenter2018byte} apply byte-level LSTM architectures for machine comprehension, reporting higher accuracy compared to word-based models when applied to morphologically complex languages like Turkish and Russian. Similarly, \cite{zhang2015character} employ convolutional models operating on characters for classification, and found that such models can surpass their word-level counterparts in certain contexts. In addition, \citet{chung2019hierarchical} introduce hierarchical LSTM models equipped with boundary-detection mechanisms at each stage to uncover the latent hierarchical structures in text, further enhancing performance in character-level language modeling. Complementing these, \cite{kalchbrenner2016neural}'s ByteNet adopts convolutional neural layers applied to characters rather than using attention mechanisms for machine translation.

\textbf{Character-Level Transformers:} The advent of attention-based transformer architectures~\citep{vaswani2017attention} in conjunction with subword tokenization~\citep{sennrich-etal-2016-neural} has led to substantial advances in language modeling and performance on standard benchmarks. Nonetheless, the adoption of word or subword units inherently prescribes an inductive bias regarding the abstraction level at which the models process language. Seeking to bridge this with earlier encouraging findings from character-level language modeling, \citet{al2019character} employed extremely deep transformer networks, supplementing their training with auxiliary loss functions. Their approach resulted in transformer models that surpassed the prior state-of-the-art character-level LSTM language models. Despite these improvements, their models continued to lag considerably behind LLMs trained at the word level. Additionally, GPT-2~\citep{radfordgpt2} demonstrated that for extensive corpora such as the 1 billion word benchmark, byte-level language models remained less effective than their word-level counterparts.

\cite{Choe2019BridgingTG} showed that transformer-based LLMs operating at the byte level can achieve better performance than their subword-level counterparts with similar numbers of parameters. However, these byte-level models require significantly greater computational resources and extended training times. In a related vein, \cite{el2020characterbert} introduced CharFormer, a BERT variant that generates word representations via convolutional operations over character embeddings, reporting gains particularly in medical NLP tasks; nonetheless, these improvements come at the cost of higher computational expenditure. CANINE, presented by \cite{clark2022canine}, is an encoder-only model with 150M parameters that processes character sequences directly. It incorporates a deep transformer backbone—functionally akin to our own global model—and applies both a local transformer and strided convolutions to downsample input characters. This architecture achieves superior performance compared to an equivalent token-based encoder-only baseline (mBERT) on multilingual downstream tasks. ByT5~\citep{xue2022byt5} further investigated byte-level encoder-decoder models that entirely forego patching strategies. Though this model demonstrated enhanced resilience to noisy input and was able to match the performance of tokenizer-dependent models with just a quarter of the training data, the absence of patching required executing expensive attention mechanisms over all input bytes, dramatically increasing the computational burden. Switching from subword to byte-level inputs leads to a significant rise in input sequence lengths, hindering scalability. More recently, leveraging the Mamba Architecture~\citep{gu2023mamba}—which retains a fixed-size memory across extended contexts—\cite{wang2024mambabyte} trained a byte-level Mamba model, also avoiding patching. Their system surpassed byte-level transformer models under equivalent FLOP budgets at the 350M parameter scale, as measured by bits-per-byte across multiple datasets.

\textbf{Patching-based approaches:} Leveraging patching techniques has the potential to substantially reduce the elevated computational costs typically associated with byte-level LLMs, while potentially maintaining comparable levels of performance. A number of studies have reported promising outcomes on relatively modest scales of model size and training data. For example, \cite{nawrot-etal-2022-hierarchical} investigate static patching for both downsampling and upsampling within the framework of the hourglass transformer, and demonstrate superior results relative to other byte-level models at the 150M parameter scale. Building on this, \cite{nawrot-etal-2023-efficient} propose dynamic patching mechanisms, introducing methods such as an end-to-end learned boundary-predictor, a boundary-predictor trained with supervision from certain tokenizers, and an entropy-driven patching model closely related to {\textsc BLT}. Their experiments show that these methods outperform standard transformers of that period on language modeling benchmarks with 40M parameters trained over roughly 400M tokens. Additionally, \cite{lester2024training} explore the use of arithmetic coding for sequence compression to surpass the compression levels achieved by BPE, and apply an equal-info windows strategy. Although this configuration yields better performance than byte-level baselines in language modeling, it does not exceed that of subword-based baselines.

Our research is heavily influenced by, and most closely aligned with, MegaByte~\citep{yu2023megabyte}. This approach utilizes a decoder-only causal LLM framework, implementing a predetermined static patching strategy that concatenates byte representations into patches; a local model within the decoder subsequently transforms these patches back into byte sequences. The creators of MegaByte provide evidence that, at the 1B parameter scale, their method achieves comparable results to models based on tokenization on datasets comprising approximately 400B bytes. In all our experiments, we conduct ablations of MegaByte and observe that static patching does not perform as well as the leading compute-optimized tokenizer-based models when accounting for floating-point operations. We then show how {\textsc BLT} narrows this performance disparity. Similarly, \cite{slagle2024spacebyte} identify analogous shortcomings in MegaByte and propose enhancing static patching to include segmenting on whitespace and other space-like characters, as well as incorporating a local encoder module. Their work shows superior results compared to tokenized-based transformer architectures under compute-constrained settings, particularly on datasets like Github and arXiv at the 1B parameter level. In this paper, we also evaluate this model and highlight the necessity for additional architectural advances to further scale byte-level approaches and reach parity with the latest leading token-based models such as Llama 3.

\section{Limitations and Future Work}

In this study, our selection of architectural configurations is guided by training models for the optimal number of steps as established for Llama~3~\citep{dubey2024llama}. It is important to note, however, that these scaling laws were derived specifically for BPE-level transformer models and may not yield optimal ratios of data to parameter size for the {\textsc BLT} architecture. The task of developing dedicated scaling laws tailored to {\textsc BLT}, which could potentially produce even more advantageous scaling behavior for this design, is reserved for future investigation. Furthermore, the majority of our experiments have been limited to models with up to 1B parameters, and it remains plausible that architectural optima may shift when scaling up to 8B parameters or larger, thereby possibly offering improved outcomes at greater model sizes.

Current transformer codebases and libraries are optimized primarily for architectures that employ tokenizers. Although our work includes theoretical FLOP-matched comparisons and leverages efficient components like FlexAttention for non-standard layers, our implementations might still lag behind tokenizer-based counterparts regarding actual runtime performance and could be further improved through additional optimization efforts.

Whereas {\textsc BLT} relies on a distinct entropy model that is trained separately for the patching process, integrating patching within an end-to-end training framework presents a compelling avenue for subsequent exploration. In Section \ref{sec:distillation}, we report preliminary findings that suggest promising outcomes for adapting tokenizer-based models like Llama 3—which have been pre-trained on over 10T tokens—toward byte-level processing, by transferring and then freezing the pretrained global transformer weights. Continued investigation along these lines may yield techniques that preserve the strengths of bytefying while also surpassing the performance of the original tokenizer-based models, without necessitating training from the beginning.

\section{Conclusion}
\label{section:conclusion}

This work introduces the Byte Latent Transformer (\textbf{BLT}), an innovative architecture that moves beyond the standard reliance on static vocabularies in large language models. By employing a dynamic and trainable approach for assembling bytes into patches, {\textsc BLT} enables adaptive allocation of computational effort according to the complexity of the input data. This mechanism yields marked gains in both efficiency and resilience. Through an extensive scaling analysis, we show that {\textsc BLT} can achieve parity with existing tokenization-dependent architectures, such as Llama 3, for models as large as 8B parameters and up to 4T bytes of input. Additionally, {\textsc BLT} offers up to a 50\% reduction in inference FLOPs with only marginal sacrifices in evaluation performance. Notably, {\textsc BLT} facilitates concurrent scaling of both model and patch sizes under a fixed inference compute budget, a characteristic particularly beneficial in practical compute-constrained scenarios. By operating directly on byte sequences, {\textsc BLT} enhances the model’s capability to manage infrequent or anomalous data, conferring greater robustness to noisy samples and deepening the representation of sub-word structures. Collectively, these findings establish {\textsc BLT} as a strong candidate to replace traditional tokenization approaches, offering an adaptable, resource-efficient, and robust framework for future large-scale language modeling.

\end{document}